% GENERATED FILE. Edit canonical lessons, then run npm run generate:books.
% canonical-source-hash: fnv1a64:529111e3430158cc
% canonical-lessons: GU-C225-cakkar, GU-C225-ulti-thavi, GU-C225-pacvu, GU-C225-bagasu-khavu, GU-C225-parsevo

\chapter{Not so well today}
\label{ch:gu-not-so-well-today}
\hlchaptermodality{225}

\begin{chapteropening}
\textbf{By the end of this chapter:} \emph{I can say that I feel dizzy or sick and that food does not agree with me.}

Complete the last lesson of chapter 225: I can say that I feel dizzy or sick and that food does not agree with me.
\end{chapteropening}

\section[\texorpdfstring{\gu{ચક્કર} (cakkar)}{ચક્કર (cakkar)}]{\gu{ચક્કર} (cakkar) — dizziness, a turn}

\label{lesson:GU-C225-cakkar}



\begin{quote}
Before the new one: say the Gujarati for to throb, then the Gujarati for breath.
\end{quote}

\subsection*{You'll want to know: \gu{ચક્કર}}
\textbf{\gu{ચક્કર}} — \emph{cakkar} — \textquotedblleft{}dizziness, a turn\textquotedblright{}.

\begin{quote}
\textbf{\gu{શું} \gu{તમને} \gu{ચક્કર} \gu{આવે} \gu{છે}?} — \emph{śũ tamne cakkar āve che?} — \textquotedblleft{}Do you feel dizzy?\textquotedblright{}
\end{quote}

A neuter noun. Word for word, does a turn come to you.

\begin{cousinweb}
\textbf{\gu{ચક્કર}} (\emph{cakkar}) is usually traced to Sanskrit \emph{cakra}, \textquotedblleft{}wheel\textquotedblright{}, so dizziness is the head going round. CDIAL lists only the inherited Gujarati \emph{cāk}, \textquotedblleft{}wheel\textquotedblright{}, under \emph{cakra}, not \emph{cakkar} itself (CDIAL 4538).
\end{cousinweb}

\begin{grammarlens}[title={What you've built}]
One more word for everyday life at A2.
\end{grammarlens}

\subsection*{Your turn}
\begin{itemize}
\raggedright
  \item \emph{Say it:} \emph{cakkar}
  \item \emph{Say it:} \emph{cakkar}, once more
  \item \emph{Say it:} \emph{śũ tamne cakkar āve che?} — Do you feel dizzy?
  \item \emph{Recall:} say the Gujarati for to throb, then the Gujarati for breath, then say \emph{cakkar} again
\end{itemize}

\begin{tcolorbox}[breakable,colback=teal!4,colframe=teal!35!black,title={Before you move on}]
What does \gu{ચક્કર} mean? (\textquotedblleft{}Dizziness, a turn\textquotedblright{}.) Say it once more.
\end{tcolorbox}
\section[\texorpdfstring{\gu{ઊલટી} \gu{થવી} (ūlṭī thavī)}{ઊલટી થવી (ūlṭī thavī)}]{\gu{ઊલટી} \gu{થવી} (ūlṭī thavī) — to vomit, to be sick}

\label{lesson:GU-C225-ulti-thavi}



\begin{quote}
Before the new one: say the Gujarati for breath, then the Gujarati for dizziness.
\end{quote}

\subsection*{You'll want to know: \gu{ઊલટી} \gu{થવી}}
\textbf{\gu{ઊલટી} \gu{થવી}} — \emph{ūlṭī thavī} — \textquotedblleft{}to vomit, to be sick\textquotedblright{}.

\begin{quote}
\textbf{\gu{શું} \gu{તમને} \gu{ઊલટી} \gu{થાય} \gu{છે}?} — \emph{śũ tamne ūlṭī thāy che?} — \textquotedblleft{}Are you being sick?\textquotedblright{}
\end{quote}

Word for word, vomiting happens: \textbf{\gu{ઊલટી}}, \emph{ūlṭī}, and \textbf{\gu{થવું}}, \emph{thavũ}.

\begin{cousinweb}
\textbf{\gu{ઊલટી} \gu{થવી}} (\emph{ūlṭī thavī}) is built from \emph{ūlṭī}, \textquotedblleft{}vomiting\textquotedblright{}, the noun of \emph{ūlṭũ}, \textquotedblleft{}turned over\textquotedblright{} (EtymDB), and \emph{thavũ}, \textquotedblleft{}to become\textquotedblright{}. CDIAL traces the family to \emph{*ullaṭati}, \textquotedblleft{}turns over\textquotedblright{}, a reconstructed form (CDIAL 2368).
\end{cousinweb}

\begin{grammarlens}[title={What you've built}]
One more word for everyday life at A2.
\end{grammarlens}

\subsection*{Your turn}
\begin{itemize}
\raggedright
  \item \emph{Say it:} \emph{ūlṭī thavī}
  \item \emph{Say it:} \emph{ūlṭī thavī}, once more
  \item \emph{Say it:} \emph{śũ tamne ūlṭī thāy che?} — Are you being sick?
  \item \emph{Recall:} say the Gujarati for breath, then the Gujarati for dizziness, then say \emph{ūlṭī thavī} again
\end{itemize}

\begin{tcolorbox}[breakable,colback=teal!4,colframe=teal!35!black,title={Before you move on}]
What does \gu{ઊલટી} \gu{થવી} mean? (\textquotedblleft{}To vomit, to be sick\textquotedblright{}.) Say it once more.
\end{tcolorbox}
\section[\texorpdfstring{\gu{પચવું} (pacvũ)}{પચવું (pacvũ)}]{\gu{પચવું} (pacvũ) — to digest, to be digested}

\label{lesson:GU-C225-pacvu}



\begin{quote}
Before the new one: say the Gujarati for dizziness, then the Gujarati for to vomit.
\end{quote}

\subsection*{You'll want to know: \gu{પચવું}}
\textbf{\gu{પચવું}} — \emph{pacvũ} — \textquotedblleft{}to digest, to be digested\textquotedblright{}.

\begin{quote}
\textbf{\gu{મને} \gu{દૂધ} \gu{પચતું} \gu{નથી}.} — \emph{mane dūdh pactũ nathī} — \textquotedblleft{}Milk does not agree with me.\textquotedblright{}
\end{quote}

Word for word, milk is not digested to me.

\begin{cousinweb}
\textbf{\gu{પચવું}} (\emph{pacvũ}) is a tadbhava, a word inherited from Sanskrit and reshaped by sound change. It comes from \emph{pacyate}, \textquotedblleft{}is cooked, is digested\textquotedblright{} (CDIAL 7654).
\end{cousinweb}

\begin{grammarlens}[title={What you've built}]
One more word for everyday life at A2.
\end{grammarlens}

\subsection*{Your turn}
\begin{itemize}
\raggedright
  \item \emph{Say it:} \emph{pacvũ}
  \item \emph{Say it:} \emph{pacvũ}, once more
  \item \emph{Say it:} \emph{mane dūdh pactũ nathī} — Milk does not agree with me.
  \item \emph{Recall:} say the Gujarati for dizziness, then the Gujarati for to vomit, then say \emph{pacvũ} again
\end{itemize}

\begin{tcolorbox}[breakable,colback=teal!4,colframe=teal!35!black,title={Before you move on}]
What does \gu{પચવું} mean? (\textquotedblleft{}To digest, to be digested\textquotedblright{}.) Say it once more.
\end{tcolorbox}
\section[\texorpdfstring{\gu{બગાસું} \gu{ખાવું} (bagāsũ khāvũ)}{બગાસું ખાવું (bagāsũ khāvũ)}]{\gu{બગાસું} \gu{ખાવું} (bagāsũ khāvũ) — to yawn}

\label{lesson:GU-C225-bagasu-khavu}



\begin{quote}
Before the new one: say the Gujarati for to vomit, then the Gujarati for to digest.
\end{quote}

\subsection*{You'll want to know: \gu{બગાસું} \gu{ખાવું}}
\textbf{\gu{બગાસું} \gu{ખાવું}} — \emph{bagāsũ khāvũ} — \textquotedblleft{}to yawn\textquotedblright{}.

\begin{quote}
\textbf{\gu{શું} \gu{બાળક} \gu{બગાસું} \gu{ખાય} \gu{છે}?} — \emph{śũ bāḷak bagāsũ khāy che?} — \textquotedblleft{}Is the child yawning?\textquotedblright{}
\end{quote}

Word for word, to eat a yawn: \textbf{\gu{બગાસું}}, \emph{bagāsũ}, and \textbf{\gu{ખાવું}}, \emph{khāvũ}.

\begin{cousinweb}
\textbf{\gu{બગાસું} \gu{ખાવું}} (\emph{bagāsũ khāvũ}) is built from \emph{bagāsũ}, \textquotedblleft{}a yawn\textquotedblright{}, and the inherited verb \emph{khāvũ}, \textquotedblleft{}to eat\textquotedblright{}. The origin of \emph{bagāsũ} is not settled: CDIAL lists no Gujarati form for it.
\end{cousinweb}

\begin{grammarlens}[title={What you've built}]
One more word for everyday life at A2.
\end{grammarlens}

\subsection*{Your turn}
\begin{itemize}
\raggedright
  \item \emph{Say it:} \emph{bagāsũ khāvũ}
  \item \emph{Say it:} \emph{bagāsũ khāvũ}, once more
  \item \emph{Say it:} \emph{śũ bāḷak bagāsũ khāy che?} — Is the child yawning?
  \item \emph{Recall:} say the Gujarati for to vomit, then the Gujarati for to digest, then say \emph{bagāsũ khāvũ} again
\end{itemize}

\begin{tcolorbox}[breakable,colback=teal!4,colframe=teal!35!black,title={Before you move on}]
What does \gu{બગાસું} \gu{ખાવું} mean? (\textquotedblleft{}To yawn\textquotedblright{}.) Say it once more.
\end{tcolorbox}
\section[\texorpdfstring{\gu{પરસેવો} (parsevo)}{પરસેવો (parsevo)}]{\gu{પરસેવો} (parsevo) — sweat}

\label{lesson:GU-C225-parsevo}



\begin{quote}
Before the new one: say the Gujarati for to digest, then the Gujarati for to yawn.
\end{quote}

\subsection*{You'll want to know: \gu{પરસેવો}}
\textbf{\gu{પરસેવો}} — \emph{parsevo} — \textquotedblleft{}sweat\textquotedblright{}.

\begin{quote}
\textbf{\gu{શું} \gu{તમને} \gu{પરસેવો} \gu{થાય} \gu{છે}?} — \emph{śũ tamne parsevo thāy che?} — \textquotedblleft{}Are you sweating?\textquotedblright{}
\end{quote}

A masculine noun.

\begin{cousinweb}
\textbf{\gu{પરસેવો}} (\emph{parsevo}) is a tadbhava, a word inherited from Sanskrit and reshaped by sound change. It comes from \emph{prasveda}, \textquotedblleft{}great sweat\textquotedblright{} (CDIAL 8898).
\end{cousinweb}

\begin{grammarlens}[title={What you've built}]
One more word for everyday life at A2.
\end{grammarlens}

\subsection*{Your turn}
\begin{itemize}
\raggedright
  \item \emph{Say it:} \emph{parsevo}
  \item \emph{Say it:} \emph{parsevo}, once more
  \item \emph{Say it:} \emph{śũ tamne parsevo thāy che?} — Are you sweating?
  \item \emph{Recall:} say the Gujarati for to digest, then the Gujarati for to yawn, then say \emph{parsevo} again
\end{itemize}

\begin{tcolorbox}[breakable,colback=teal!4,colframe=teal!35!black,title={Before you move on}]
What does \gu{પરસેવો} mean? (\textquotedblleft{}Sweat\textquotedblright{}.) Say it once more.
\end{tcolorbox}
